\section{Dynamical deconcentration of critical edges at an exact index}
\label{sec:v45-critical-deconcentration}
We now use the arithmetic local law to estimate the actual edge
coefficients, rather than only reconstruct their sum. The argument
compares a positive endpoint collar with a positive endpoint-selected
local probability. No count of words or separation of critical values
enters the comparison. The final conclusions are pointwise statements
at the original exact displacement and return index.

\subsection{A two-endpoint small-source estimate}
Write $p(x)$ for the tangential collision velocity, and for a fixed
bounded interval $J$ put
\begin{align}
 E_{n,k,m,R}(t;J)&=\{K_{n,R}=k,\ N_{n,R}=m,\ T_{n,R}-t\in J\},
                                                        \notag\\
 E^{(d)}_{n,k,m,R}(t;J)&=E_{n,k,m,R}(t;J)\cap
       \{|p(x)|\le d,\ |p((F_R^*)^n x)|\le d\}.
                                                        \label{eq:v45-strip-event}
\end{align}
The terminal state is $T_R^m x$ on this unchanged exact event.

\begin{lemma}[Local mass with two normal endpoint strips]
\label{lem:v45-two-strip-local}
There is $C<\infty$, independent of $0<d<1/4$ and $J$, such that
\begin{equation}\label{eq:v45-two-strip-local}
 \limsup_{m\to\infty}\sup_{\substack{R\in I,\,1\le n\le m\\
                                     k\in\Z^2,\,t\in\R}}
       m^2\nu_R^*(E^{(d)}_{n,k,m,R}(t;J))
                       \le C d^2|J|.
\end{equation}
The interval and strip width are fixed before the collision count tends
to infinity. The assertion is an upper bound and does not replace the
arithmetic factor by one or assert a positive denominator.
\end{lemma}
\begin{proof}
Choose a smooth $0\le b_d\le1$ on the collision cylinder, equal to
one for $|p|\le d$ and zero for $|p|\ge2d$. Put $a_{R,d}=\eta_Rb_d$.
The initial vector $a_{R,d}\nu$ and terminal multiplier $a_{R,d}$ are
bounded in the relevant collision norms for each fixed $d$, uniformly
in $R$. Their physical masses satisfy
\begin{equation}\label{eq:v45-small-endpoint-mass}
 \nu(a_{R,d})\le\nu\{|p|<2d\}=2d.
\end{equation}
No estimate uniform in $d$ for their strong norms is needed.

Repeat the finite-branch fixed-band inverse of
Theorem~\ref{thm:v41-uniform-transition} with these two endpoints.
The spectral charts, eigenvalues, real maxima and covariance bounds
are unchanged; only the scalar amplitudes become
\[
 A_{j,d}(R,\theta)=
     \ell(M_{a_{R,d}}\Pi_j(R,\theta)(a_{R,d}\nu)).
\]
For fixed $d$ these amplitudes and their frequency derivatives have
uniform bounds and are continuous in $R$. To check the latter use
finite-time physical pairings with the two factors $a_{R,d}$, as in
Lemma~\ref{lem:v41-scalar-continuity}. Finite itineraries and section
membership persist off a null set at the limiting radius. Normalized
powers converge uniformly to the scalar projection on the local
contour. Thus the same uniform-limit argument applies to these
endpoints, without asserting strong-norm continuity of $\eta_R$.

Take any parameter and target sequence in the supremum and pass to
$R_m\to R_0$ inside the finite chart cover. A peak branch whose
limiting maximum has modulus less than one contributes zero in the
limit. Every remaining limiting peak is an actual resonance. The
physical projection formula there gives
\begin{equation}\label{eq:v45-strip-residue-bound}
 |A_{j,d}(R_0,a_j(R_0))|
   =|\nu(a_{R_0,d}\overline q_j)\nu(a_{R_0,d}q_j)|
       \le4d^2,
\end{equation}
where $|q_j|=1$ almost everywhere. This uses all resonances, including
those reached through parameter transitions. Each damping factor has
modulus at most one. The complex Gaussians have a common bound from
$\operatorname{Re}H_j\ge aI_4$, for every target argument. The number
of peak charts is fixed independently of the endpoints and of the
outer roof band. Indeed every exact roof resonance has roof coordinate
zero, as specified in Lemma~\ref{lem:v41-moving-maxima}.

Let $q_B$ be a positive Fourier majorant of the interval indicator,
with mass $|J|+2\pi/B$. For each fixed $B$, the band inverse and the
preceding peak estimate bound the limsup of its endpoint-selected
expectation, multiplied by $m^2$, by $Cd^2(|J|+2\pi/B)$.
In this step the roof coordinates of surviving limiting peaks are
zero, so their test factor is $\widehat q_B(0)=\int q_B$.
The positive exact source permits the envelope inequality before any
limit is taken. First let $m\to\infty$, then $B\to\infty$.
Equation~\eqref{eq:v45-strip-event} is bounded by that selected source,
which proves \eqref{eq:v45-two-strip-local}. The subsequence argument
proves uniformity even when the arithmetic type changes with $R$.
\end{proof}

\subsection{Coalescing critical values}
Let $\mathcal C^\varepsilon_{n,k,m,R}$ be all $\varepsilon$-protected
normal-to-normal critical words whose original source has
$N_{n,R}=m$ and $K_{n,R}=k$. This is a finite set for each $m$.
The coefficient $J_z$ is that in \eqref{eq:v45-intrinsic-jump}.
Define the positive atomic measure
\begin{equation}\label{eq:v45-critical-coefficient-measure}
 \mathcal J^\varepsilon_{n,k,m,R}
                    =\sum_{z\in\mathcal C^\varepsilon_{n,k,m,R}}
                                              J_z\delta_{t_z}.
\end{equation}
Equal critical values are combined by addition in this measure.
The definition does not change any of their intrinsic coefficients.

\begin{theorem}[A local bound for the complete regular critical cluster]
\label{thm:v45-critical-cluster}
There are $C<\infty$ and $h_*>0$, independent of $\varepsilon$,
such that, with $h_\varepsilon=h_*\varepsilon^6$, for every fixed
$0<h<h_\varepsilon$,
\begin{equation}\label{eq:v45-cluster-bound}
 \limsup_{m\to\infty}\sup_{R,n,k,t}
    m^2\mathcal J^\varepsilon_{n,k,m,R}([t-h,t+h])
                                      \le C h.
\end{equation}
Here and below $1\le n\le m$. In particular,
\begin{equation}\label{eq:v45-atom-smallness}
 \sup_{R,n,k,t}m^2\mathcal J^\varepsilon_{n,k,m,R}(\{t\})
                              \longrightarrow0.
\end{equation}
No separation of critical values, bound on the number of words, or
triviality of a section residue is assumed.
\end{theorem}
\begin{proof}
For each critical value in $[t-h,t+h]$, use the actual source collar
$U_z(h)$ of Proposition~\ref{prop:v45-mass-collar}. These collars are
disjoint and keep exactly the prescribed labels $(n,k,m)$. Their roof
values lie in $[t-h,t+2h]$. Both endpoint tangential velocities have
modulus at most $A_*\sqrt h$. Thus
\begin{equation}\label{eq:v45-positive-collar-comparison}
 a_* h\,
       \mathcal J^\varepsilon_{n,k,m,R}([t-h,t+h])
 \le\nu_R^*\bigl(E^{(A_*\sqrt h)}_{n,k,m,R}
                                  (t;[-h,2h])\bigr).
\end{equation}
Reduce $h_\varepsilon$ so that the strip width is below $1/4$.
For this fixed $h$, Lemma~\ref{lem:v45-two-strip-local} bounds the
limsup of the right side times $m^2$ by $C h^2$.
Division by $a_* h$ proves \eqref{eq:v45-cluster-bound}.
An atom is bounded by every interval in that estimate. Let $h$ tend
to zero after the limsup in $m$ to obtain
\eqref{eq:v45-atom-smallness}. The auxiliary roof interval is used
only to bound a positive source; neither exact integer label is
averaged, and the conclusion concerns the original individual roof
value, including coincident critical values.
\end{proof}

\subsection{A pointwise-small part of the original source}
For any sequence $r_m\downarrow0$, eventually less than
$h_\varepsilon$, restrict the original exact-label source to
\[
 \mathcal U^\varepsilon_{n,k,m,R}(r_m)
       =\bigcup_{z\in\mathcal C^\varepsilon_{n,k,m,R}}U_z(r_m).
\]
Write $p^{\varepsilon,\mathrm{crit}}_{n,k,m,R}$ for its roof density.
The source has not been replaced by a probability with a new
normalization; it is a positive restriction of $\nu_R^*$.

\begin{theorem}[Uniform pointwise removal of regular critical collars]
\label{thm:v45-critical-source-smallness}
For every fixed $\varepsilon>0$ and every $r_m\downarrow0$,
\begin{equation}\label{eq:v45-critical-density-smallness}
 \sup_{R,n,k}m^2\|p^{\varepsilon,\mathrm{crit}}_{n,k,m,R}\|_\infty
                                        \longrightarrow0.
\end{equation}
For arbitrary measurable insertions of modulus at most $M$, the same
restricted-source conclusion holds with the upper bound multiplied
by $M$. It is uniform over this whole bounded insertion class.

There is also a coefficient-level version. For each $z$ choose a
coefficient $\beta_z$ with $|\beta_z|\le M J_z$ and a function
$0\le\vartheta_z\le1$ supported in $(0,r_m)$. Then
\begin{equation}\label{eq:v45-jump-profile-smallness}
 j^\varepsilon_{n,k,m,R}(t)
       =\sum_{z\in\mathcal C^\varepsilon_{n,k,m,R}}
                       \beta_z\vartheta_z(t-t_z)
 \quad\hbox{satisfies}\quad
 \sup_{R,n,k}m^2\|j^\varepsilon_{n,k,m,R}\|_\infty\longrightarrow0.
\end{equation}
In particular this applies to localized step profiles with their
original intrinsic regular critical coefficients.
\end{theorem}
\begin{proof}
The collar density upper bound in
\eqref{eq:v45-collar-density}, disjointness, and positivity give
\begin{equation}\label{eq:v45-critical-density-cluster}
 p^{\varepsilon,\mathrm{crit}}_{n,k,m,R}(t)
       \le A_*
              \mathcal J^\varepsilon_{n,k,m,R}([t-r_m,t]).
\end{equation}
Fix $h>0$ in \eqref{eq:v45-cluster-bound}. Eventually $r_m<h$, so
that estimate bounds the limsup of the left side's essential-supremum
norm, multiplied by $m^2$, by $C h$.
Now let $h\downarrow0$. This proves
\eqref{eq:v45-critical-density-smallness} for every such sequence;
no spectral constant is evaluated at $r_m^{-1}$.
A bounded insertion is dominated in variation by $M$ times the same
positive source, so its density has the same pointwise domination.
For the coefficient version use directly
$|j^\varepsilon(t)|\le M\mathcal J^\varepsilon([t-r_m,t])$.
The identical argument proves \eqref{eq:v45-jump-profile-smallness},
even when many supports overlap or several critical values coincide.
\end{proof}

\begin{corollary}[Removal from the complete signed correction]
\label{cor:v45-critical-unsmoothing}
Let $K_B$ be the unchanged convolution kernel in
Theorem~\ref{thm:v43-arithmetic-raw-criterion}. For either the weighted
critical-source density or the step profile in the preceding theorem,
write $f^{\mathrm{crit}}$ for that component. Then
\begin{equation}\label{eq:v45-critical-all-band}
 \sup_{B>0}\sup_{R,n,k}m^2
             \|f^{\mathrm{crit}}-K_B*f^{\mathrm{crit}}\|_\infty
                                         \longrightarrow0.
\end{equation}
For the positive-source decomposition $p=p^{\varepsilon,\mathrm{crit}}
+p^{\varepsilon,\mathrm{rest}}$ this gives, in the same uniform
pointwise sense,
\begin{equation}\label{eq:v45-remaining-source}
 \mathscr D_{B,n,R}
   =p^{\varepsilon,\mathrm{rest}}-K_B*p^{\varepsilon,\mathrm{rest}}
                                      +o(m^{-2}).
\end{equation}
The coefficient and source normalizations in this identity are the
original ones. Its error estimate is uniform in $B>0$.
\end{corollary}
\begin{proof}
The scale-independent identity $\|K_B\|_1=\|K_1\|_1$ gives
$\|f-K_B*f\|_\infty\le(1+\|K_1\|_1)\|f\|_\infty$.
Apply Theorem~\ref{thm:v45-critical-source-smallness} and use linearity.
The uniformity in $B$ here is a convolution inequality for a part of
the source already proved pointwise small. It is not an estimate for
transfer-operator constants on growing frequency bands.
\end{proof}

\paragraph{Boundary of the estimate.}
The argument estimates the entire regular critical cluster under the
fixed exponentially relaxed margin envelope, rather than multiplying an individual edge
bound by an uncontrolled word count. The margin is fixed in
\eqref{eq:v45-cluster-bound}; no usable dependence of the asymptotic
threshold on a margin tending to zero is asserted. Critical words
violating this envelope near grazing or section decisions, other boundary contributions,
and the remaining noncritical signed inverse are not covered by
\eqref{eq:v45-remaining-source}. The full raw target retains
$\mathcal L_{m,R}$, and at fixed radius
$c\mathfrak a_Rg_{\Omega_R}$, as its main term. The pointwise-small
positive restriction is returned to the exact source identity; it
is not a replacement of the requested singleton by a roof window.
